\chapter{Beyond Gaussianity: Laplace Innovations and the Gaussian-Message
Approximation}\label{ch:laplace}\label{sec:bp-nongaussian}
% The full title wraps the running head onto a second line, which no
% headheight can absorb gracefully; the mark carries a one-line form.
\chaptermark{Beyond Gaussianity: Laplace Innovations}

Every exact result of Chapter~\ref{ch:gaussian} rests on the closure of
the Gaussian family under the two operations the model performs,
multiplying densities and marginalising. This chapter removes that
property in the most controlled way available and answers RQ3. We keep
the AR(1) recursion, the OU channel, and the factor-graph structure, and
change only the innovation law from Gaussian to Laplace: heavy-tailed and
kinked at the origin, yet log-concave and analytically tractable.

\section{What changes, and what is preserved}\label{sec:l-model}

The Laplace AR(1) chain replaces the Gaussian innovations of
\eqref{eq:ar1} with Laplace ones, keeping the initial law and the
second-order structure fixed:
\begin{equation}
  a_0 \sim \Nd(0,1), \qquad
  a_k = \alpha\, a_{k-1} + \varepsilon_k, \qquad
  \varepsilon_k \sim \mathrm{Laplace}(0, b) \ \text{i.i.d.}, \qquad
  k = 1,\dots,\nsites-1 ,
  \label{eq:l-chain-prior}
\end{equation}
observed through the same per-frame OU channel. The innovation scale is
\emph{not} free: $b = \sqrt{(1-\alpha^2)/2}$, so that
$\operatorname{Var}(\varepsilon_k) = 1-\alpha^2$ and the chain has unit
marginal variance and $\operatorname{Cov}(a_j,a_k)=\alpha^{|j-k|}$, exactly
as in the Gaussian case. The Gaussian and Laplace chains therefore agree on
every second moment and differ \emph{only} beyond them.

The score--posterior
identity \eqref{eq:tweedie-joint} holds for any $p_0$. The posterior
factorises over the same tree-shaped factor graph, so the sum--product recursion
\eqref{eq:bp-fwd}--\eqref{eq:bp-bwd} still computes exact smoothing marginals
in one forward and one backward sweep. And the OU likelihood factors are still
Gaussian in $a_k$. What is lost is that the messages are no longer
Gaussian
and the score ceases to be linear in $x$.

\section{Where the score stops being linear}
\label{sec:l-nonlinear}\label{sec:l-K1}

The score--posterior identity \eqref{eq:tweedie-joint} is indifferent to the
prior. Componentwise, for the chain \eqref{eq:l-chain-prior} under the OU
channel, with $\mu := e^{-t}$,
\begin{equation}
  \score_k(x,t) \;=\; \frac{\mu\, \E[a_k \mid x] \;-\; x_k}{\Deltat} .
  \label{eq:l-score-identity}
\end{equation}
For Gaussian innovations $\E[a_k \mid x]$ is a linear function of
$x$ and \eqref{eq:l-score-identity} is the linear score $-\Qt x$; for Laplace
innovations the posterior mean is nonlinear, and the score inherits the
nonlinearity exactly.

The two chains agree on every second moment by construction
(Section~\ref{sec:l-model}); they differ first at fourth order. The
\emph{innovation} carries excess kurtosis $3$: for a centred Laplace variable
of variance $v$, $\kappa_4 = 3v^2$, so the standardised fourth cumulant is
$3$, against $0$ for a Gaussian innovation.

The chain \emph{state} does not inherit that number unchanged, because
$a_k$ is a weighted sum of one Gaussian seed and $k$ independent Laplace
innovations, and cumulants of independent sums add with the fourth power of
the coefficients. From \eqref{eq:l-chain-prior},
$a_k = \alpha^k a_0 + \sum_{j=1}^{k}\alpha^{k-j}\varepsilon_j$ with
$\kappa_4(a_0) = 0$ and
$\kappa_4(\varepsilon_j) = 3(1-\alpha^2)^2$, so
\begin{equation}
    \kappa_4(a_k)
    = 3(1-\alpha^2)^2\sum_{m=0}^{k-1}\alpha^{4m}
    = 3\,\frac{1-\alpha^2}{1+\alpha^2}\,\bigl(1-\alpha^{4k}\bigr),
    \label{eq:l-state-kurtosis}
\end{equation}
which is also the excess kurtosis of $a_k$ since
$\operatorname{Var}(a_k) = 1$. Two consequences matter later. The state's
non-Gaussianity is much milder than the innovation's --- at $\alpha = 0.85$
the bulk value is $3 \times 0.2775/1.7225 \approx 0.48$, not $3$ --- and it
depends on the position $k$: the chain is covariance-stationary by
construction but \emph{not} strictly stationary, because the Gaussian initial
law is not the stationary law of the Laplace recursion, and the shape settles
towards its bulk value \eqref{eq:l-state-kurtosis} only as $\alpha^{4k}$
dies away.

Under the variance-preserving forward process, fourth cumulants of the
independent sum $x_{t,k} = e^{-t}a_k + \sqrt{\Deltat}\,z_k$ scale as the
fourth power of the coefficients, and $\kappa_4(z) = 0$, so
\begin{equation}
    \boxed{\;\kappa_4(x_{t,k}) = e^{-4t}\,\kappa_4(a_k)\;}
    \label{eq:l-noised-kurtosis}
\end{equation}
with $\kappa_4(a_k)$ as in \eqref{eq:l-state-kurtosis}. Hence Laplace-driven
non-Gaussianity survives every finite diffusion time but decays as $e^{-4t}$,
while the Gaussian case remains Gaussian throughout.

Figure~\ref{fig:l-nonlinear} shows how this appears in the score. 

\begin{figure}[t]\centering
\includegraphics[width=\linewidth]{fig_laplace_nonlinear.pdf}
\caption[Laplace innovations: same variance, nonlinear score]{What changes
when the innovations go Laplace ($\alpha = 0.85$, so
$b = \sqrt{(1-\alpha^2)/2}$). Vermilion and solid is Laplace throughout; blue
and dashed is Gaussian. \emph{Left:} the two innovation laws on a log scale.
They carry the same variance $1-\alpha^2$ by construction and differ only in
shape --- the kink at the origin and the exponential rather than Gaussian
tail --- which the fourth cumulant summarises as excess kurtosis $+3$ against
$0$. \emph{Centre and right:} the joint score \eqref{eq:l-score-identity} at
the middle frame of a chain of $\nsites = 9$, as the observation $x_k$ at that
frame is swept with the others held at zero, computed by grid BP
(Section~\ref{sec:l-grid}) against the covariance-matched Gaussian chain ---
which is exactly the estimator the Gaussian closure returns
(Theorem~\ref{thm:l-lmmse}). At $t = 0.05$ the Laplace score saturates near
$\pm 5.2$ while the Gaussian one grows without bound; at $t = 1$, on the same
vertical scale, the two coincide.}
\label{fig:l-nonlinear}
\end{figure}

\section{Where closed forms stop}\label{sec:l-chain}

The tree structure is untouched by any of this, so the message decomposition
survives in full.

\begin{theorem}[Joint score from chain messages]\label{thm:l-markov-score}
For any chain prior, in particular \eqref{eq:l-chain-prior}, and the OU
likelihood, the joint score is \eqref{eq:l-score-identity} with
\begin{equation}
  \E[a_k \mid x]
  = \frac{\displaystyle\int a_k\; \lmsg_k(a_k)\;
      \psi^{(k)}_{\mathrm{ob}}(a_k)\; \rmsg_k(a_k)\;\dd a_k}
     {\displaystyle\int \lmsg_k(a_k)\;
      \psi^{(k)}_{\mathrm{ob}}(a_k)\; \rmsg_k(a_k)\;\dd a_k},
  \label{eq:l-markov-score}
\end{equation}
the messages being those of Proposition~\ref{prop:bp-convA}.
\end{theorem}

Computational locality survives. What does not survive is the collapse of message updates into two parameters.

\begin{remark}[Where analytic closure breaks]\label{rmk:l-breakdown}
The first message can be evaluated explicitly. With
\[
M(a'\mid a)
=
\frac{1}{2b}
\exp\!\left(-\frac{|a'-\alpha a|}{b}\right),
\qquad
p(x'\mid a')
=
\Nd(x';\mu a',\Deltat),
\]
set $r=a'-\alpha a$ and $y=x'-\mu\alpha a$. Then
\begin{align}
I(y)
&:=
\int_{\R} M(a'\mid a)\,
        \Nd(x';\mu a',\Deltat)\,\dd a' \\
&=
\frac{1}{2b\sqrt{2\pi\Deltat}}
\int_{\R}
\exp\!\left[
-\frac{|r|}{b}
-\frac{(y-\mu r)^2}{2\Deltat}
\right]\dd r .
\end{align}
Splitting at the single kink $r=0$ and completing the square on each
half-line gives
\begin{equation}
\boxed{
I(y)
=
\frac{e^{\Deltat/(2\mu^2b^2)}}{2\mu b}
\left[
e^{-y/(\mu b)}
\Phi\!\left(
\frac{y}{\sqrt{\Deltat}}
-\frac{\sqrt{\Deltat}}{\mu b}
\right)
+
e^{y/(\mu b)}
\Phi\!\left(
-\frac{y}{\sqrt{\Deltat}}
-\frac{\sqrt{\Deltat}}{\mu b}
\right)
\right],
}
\end{equation}
for $\mu>0$, where $\Phi$ is the standard Gaussian cdf. Thus the first
Laplace--Gaussian convolution closes in elementary exponentials and error
functions.

Closure is already lost at the next update. The incoming message is now a
sum of terms of the form
\[
e^{ca}\Phi(da+e),
\]
so multiplication by the Gaussian likelihood and the next Laplace transition
requires integrals such as
\[
\int_{-\infty}^{\alpha a'}
e^{-q a^2+\ell a}\Phi(da+e)\,\dd a
\qquad\text{and}\qquad
\int_{\alpha a'}^\infty
e^{-q a^2+\ell a}\Phi(da+e)\,\dd a .
\]
These involve Gaussian cdfs inside truncated Gaussian integrals---equivalently,
higher-dimensional Gaussian probabilities---and no longer belong to the
one-dimensional normal--Laplace family. Each propagation step therefore
increases the functional complexity, unlike the Gaussian chain, where
Gaussian-times-Gaussian integration preserves the same two-parameter family
indefinitely.

Changing coordinates does not remove the obstruction. Define the innovations
\[
r_0=a_0,
\qquad
r_k=a_k-\alpha a_{k-1},
\]
so that
\[
a_k=\sum_{j=0}^k \alpha^{k-j}r_j,
\qquad a=Br,
\qquad
B_{kj}=\alpha^{k-j}\mathbf 1_{\{k\ge j\}}.
\]
The prior is then factorised, but the Gaussian likelihood contains
\[
r^\T B^\T B r,
\]
with
\[
(B^\T B)_{ij}
=
\alpha^{|i-j|}
\frac{1-\alpha^{2(L-\max\{i,j\})}}
     {1-\alpha^2},
\]
which is dense whenever $\alpha\neq0$. The dependence has therefore only
moved from the prior to the likelihood. Exact BP remains valid, but its
messages must be represented numerically; the next section describes that
representation.
\end{remark}

\section{Grid belief propagation: the discretised recursion}
\label{sec:l-grid}

\subsection{The representation and the update}

A message is a density on $\R$. We represent it by its values on $\ngrid$
equispaced points of a truncated interval,
\begin{equation}
  g_i \;=\; -A + i\,h, \qquad h \;=\; \frac{2A}{\ngrid - 1},
  \qquad i = 0, \dots, \ngrid-1,
  \label{eq:l-grid}
\end{equation}
so that each message is a vector in $\R^{\ngrid}$ and every integral in
\eqref{eq:bp-fwd}--\eqref{eq:bp-bwd} becomes a composite trapezoidal
quadrature with weights
\begin{equation}
  w_i \;=\;
  \begin{cases}
    h/2, & i \in \{0,\ \ngrid-1\},\\
    h,   & \text{otherwise},
  \end{cases}
  \label{eq:l-trap-weights}
\end{equation}
the halved endpoints being what distinguishes the trapezoidal from the
rectangle rule (numerically immaterial here, since the truncation level of
$A = 8$ keeps the endpoint values below $10^{-9}$, but the implementation
carries them and the equations should say what the implementation does). Two arrays carry the
model. The transition enters as a matrix, the likelihood as a diagonal:
\begin{equation}
  K_{ij} \;=\; M\big(g_i \mid g_j\big)
        \;=\; \frac{1}{2b}\, e^{-|g_i - \alpha g_j| / b},
  \qquad
  \big(\ell_k\big)_i
        \;=\; \Nd\big(x_k;\, \mu g_i,\, \Deltat\big).
  \label{eq:l-grid-operators}
\end{equation}
The forward and backward sweeps are then the trapezoidal quadratures of
Proposition~\ref{prop:bp-convA}, with $\odot$ the elementwise product and
$w = (w_i)$ the weight vector of \eqref{eq:l-trap-weights}:
\begin{align}
  f_0 &\;\propto\; p_0 \odot \ell_0,
  &f_k &\;\propto\; \ell_k \odot K\big( w \odot f_{k-1} \big),
  \quad k = 1,\dots,\nsites-1,
  \label{eq:l-grid-fwd}\\[2pt]
  \beta_{\nsites-1} &\;=\; \mathbf{1},
  &\beta_k &\;\propto\; K^\T \big( w \odot \ell_{k+1} \odot \beta_{k+1} \big),
  \quad k = \nsites-2,\dots,0 .
  \label{eq:l-grid-bwd}
\end{align}
The belief, posterior mean and score at frame $k$ follow from
\eqref{eq:l-score-identity}, with the same weights in the moments:
\begin{equation}
  \mathrm{bel}_k \;=\; f_k \odot \beta_k,
  \qquad
  \widehat m_k
    \;=\; \frac{\sum_i w_i\, g_i\, (\mathrm{bel}_k)_i}{\sum_i w_i\, (\mathrm{bel}_k)_i},
  \qquad
  \widehat \score_k \;=\; \frac{\mu\, \widehat m_k - x_k}{\Deltat}.
  \label{eq:l-grid-score}
\end{equation}
Each sweep step is one matrix--vector product, so the solver costs
$O(\nsites \ngrid^2)$ time and $O(\nsites \ngrid)$ space: the cost of
sum--product on a chain over a finite alphabet, with the alphabet size
replaced by the grid resolution $\ngrid$. Nothing about the algorithm has
changed; only the alphabet is now a discretisation rather than a given.

\subsection{The three choices, and what fixes them}

\paragraph{Renormalisation.} The proportionalities in
\eqref{eq:l-grid-fwd}--\eqref{eq:l-grid-bwd} are enforced after every step, the
forward messages to unit discrete mass and the backward messages to unit
maximum. This is free --- sum--product messages are defined only up to a
positive constant (Theorem~\ref{thm:tree-exact}) and the ratio
\eqref{eq:l-grid-score} is invariant --- and it is necessary: without it the
message norms drift by a factor per frame and a chain of $\nsites = 32$
underflows.

\paragraph{Truncation $A$.} The messages are beliefs about a chain of unit
marginal variance, so mass beyond $|a| = A$ is bounded by the tail the
innovations can produce. We take $A = 8$, at which the exponential tail
contributes below $10^{-9}$ of the mass for every $\alpha$ used here; the
diagnostic is the mass in the outermost grid cells, logged per run.

\paragraph{Resolution $\ngrid$.} The binding constraint is not the smoothness
of the messages but the width of the likelihood. In $a$-coordinates
$\ell_k$ has standard deviation $\sqrt{\Deltat}/\mu$, which for small $t$ is
$\sqrt{2t} + O(t^{3/2})$; if that falls below the grid spacing the likelihood
is a spike between two nodes and the quadrature ceases to see it. We therefore
require
\begin{equation}
  \sqrt{2t} \;\gtrsim\; 3h \;=\; \frac{6A}{\ngrid - 1},
  \label{eq:l-resolution}
\end{equation}
implemented in code as \codefile{src/noising.py}'s
\texttt{likelihood\_resolution\_ok} and recorded per cell by the grid
validation of \codefile{exp\_01}, which writes it as a column rather than
refusing to run: the working configuration clears it with room to spare
($3h = 0.12$ against $\sqrt{2t} = 0.20$ at the smallest reported $t$,
$t = 0.02$), so it documents the margin rather than gating on it. It is
\eqref{eq:l-resolution}, not
the innovation law, that sets the smallest diffusion time this chapter reports.

\subsection{Validation}

Grid BP is a numerical approximation of the functional recursion and can serve
as a reference only inside a validated regime, which we establish on the
exactly solvable Gaussian chain, where the closed form of
Chapter~\ref{ch:gaussian} gives the truth. For smooth, rapidly decaying
integrands --- the Gaussian validation case --- the trapezoidal rule converges
superalgebraically, so failure is abrupt rather than gradual, and the
measurement shows exactly that: at $\ngrid = 101$, $A = 8$ the relative score
error sits at the double-precision floor down to $t \approx 0.03$, has grown
to $10^{-10}$ at $t = 0.02$ and $10^{-6}$ at $t = 0.012$, and reaches the
percent scale by $t = 0.005$. The gate \eqref{eq:l-resolution} is deliberately
conservative relative to that onset --- at $\ngrid = 101$ it already demands
$t \gtrsim 0.115$ --- so every configuration this thesis reports passes it
with a wide margin rather than sitting at the failure edge. (For the Laplace
kernel the integrand has a kink along $a' = \alpha a$ and no spectral
statement is invoked; the non-Gaussian reference is certified by
self-convergence instead.) The working configuration for every
non-Gaussian experiment is $\ngrid = 401$, $A = 8$ --- gate satisfied for all
$t \geq 0.02$ with room to spare --- with a reference solution
at $\ngrid = 801$; the self-convergence error between the two remains at
least three orders of magnitude smaller than every effect measured below,
and four to five over most of the range (Figure~\ref{fig:bp-closure}, left).

Self-convergence answers whether the mesh is fine enough and cannot answer
whether the domain is wide enough, since both resolutions share the same
truncation. That question is settled separately, by the half-width sweep of
\codefile{exp\_18} ($\gdtrunccells$ cells over $A \in \{4,6,8,10\}$ and three
resolutions), which records the posterior mass reaching the outermost cells of
the domain: $\gdtruncfour$ at $A = 4$, $\gdtruncsix$ at $A = 6$,
$\gdtrunceight$ at the working $A = 8$, and $\gdtruncten$ at $A = 10$. At the
setting used here the truncated tails carry parts in $10^{6}$ of the
posterior, two orders below the smallest closure error reported in this
chapter, so the domain is not a binding approximation. Within this
regime --- resolution gated, truncation measured --- we treat grid BP as the
reference score; outside it no claim is made.

Finally, one identity governs how all of this is reported. Every score error is
an amplified posterior-mean error: for any two solvers with means $\widehat m$
and $m_{\mathrm{ref}}$, \eqref{eq:l-score-identity} gives, exactly,
\begin{equation}
  \widehat s - s_{\mathrm{ref}}
  = \frac{e^{-t}}{\Deltat}\,\big( \widehat m - m_{\mathrm{ref}} \big),
  \qquad
  \frac{e^{-t}}{\Deltat} \sim \frac{1}{2t} \;\; (t \downarrow 0).
  \label{eq:bp-master-error}
\end{equation}
Its numerical residual is logged in every experiment and stays between
$10^{-15}$ and $10^{-12}$. Equation \eqref{eq:bp-master-error} says in advance
why small diffusion times are the difficult regime: whatever an approximation
does to the posterior mean is amplified by $1/2t$ in the score.

\section{The Gaussian-message approximation, measured}
\label{sec:l-closure}

We now test \emph{Gaussian-projected BP}: the same forward--backward
recursion as grid BP, but with every transition message projected onto the
Gaussian family by moment matching. Write $\Pi_{\mathrm G} f$ for the Gaussian
density with the same mean and variance as $f/\!\int f$.

\begin{theorem}[Gaussian closure is the LMMSE estimator]
\label{thm:l-lmmse}
Consider moment-projected sum--product on the chain
\eqref{eq:l-chain-prior}, with $\alpha\neq0$, Gaussian unary likelihoods,
$\E[a_0]=0$, innovation mean zero and variance $\ivar$, and a Gaussian
initial message, or one projected by $\Pi_{\mathrm G}$. Then the posterior
means coincide exactly with those of the covariance-matched Gaussian model:
\begin{equation}
  \widehat a_{\mathrm{LMMSE}}
  =
  \operatorname{Cov}(a,x)\operatorname{Cov}(x)^{-1}x
  =
  e^{-t}\Sigz
  \bigl(e^{-2t}\Sigz+\Deltat I\bigr)^{-1}x .
  \label{eq:l-lmmse}
\end{equation}
Hence the closure returns the same estimator for every innovation law sharing
the moments $(0,\ivar)$.
\end{theorem}

\subsection{Proof of Theorem~\ref{thm:l-lmmse}}
\label{sec:lmmse-proof}

\begin{proof}
We show that every projected message is identical to the corresponding exact
message of the covariance-matched Gaussian chain.

Let $\lmsg_i$ denote the projected forward message at site $i$ and $\ell_i$
its Gaussian unary likelihood.

\emph{Base case.}
The initial message is Gaussian, either directly or after projection, with the
same first two moments as the Gaussian reference chain. The two initial
messages therefore coincide.

\emph{Forward recursion.}
Assume the incoming message is Gaussian,
\[
  \lmsg_{i-1}(a_{i-1})=\Nd(a_{i-1};m,v).
\]
Multiplication by the Gaussian likelihood $\ell_{i-1}$ remains Gaussian, with
the same precision-weighted update in both models. Propagation through
\[
  a_i=\alpha a_{i-1}+\varepsilon_i,
  \qquad
  \E[\varepsilon_i]=0,
  \qquad
  \operatorname{Var}(\varepsilon_i)=\ivar,
\]
maps its moments as
\begin{equation}
  (m,v)
  \longmapsto
  \bigl(\alpha m,\alpha^2v+\ivar\bigr).
  \label{eq:lmmse-fwd-map}
\end{equation}
Crucially, this transformation depends on the innovation distribution only
through $(0,\ivar)$. After applying $\Pi_{\mathrm G}$, the resulting Gaussian
is therefore exactly the transition message of the covariance-matched
Gaussian chain.

\emph{Backward recursion.}
The same argument applies in reverse. If the incoming Gaussian factor has
moments $(m,v)$, the backward transition integral, viewed as a density in the
preceding state, has moments
\begin{equation}
  (m,v)
  \longmapsto
  \left(\frac{m}{\alpha},
        \frac{v+\ivar}{\alpha^2}\right),
  \label{eq:lmmse-bwd-map}
\end{equation}
again depending only on the innovation variance. The condition
$\alpha\neq0$ is used here. Since the terminal message
$\rmsg_{\nsites-1}\equiv1$ is not normalisable, the induction begins at the
terminal Gaussian likelihood $\ell_{\nsites-1}$, which is the first object
actually propagated and projected.

Thus all projected forward and backward messages, and hence all posterior
beliefs, coincide with those of the covariance-matched Gaussian chain.

For jointly Gaussian zero-mean $(a,x)$,
\[
  \E[a\mid x]
  =
  \operatorname{Cov}(a,x)\operatorname{Cov}(x)^{-1}x.
\]
Under
\[
  x=e^{-t}a+\sqrt{\Deltat}\,z,
  \qquad z\perp a,
\]
we have
\begin{equation}
  \operatorname{Cov}(a,x)=e^{-t}\Sigz,
  \qquad
  \operatorname{Cov}(x)=e^{-2t}\Sigz+\Deltat I .
  \label{eq:lmmse-covs}
\end{equation}
Substitution gives \eqref{eq:l-lmmse}. Since neither
\eqref{eq:lmmse-fwd-map} nor \eqref{eq:lmmse-bwd-map} uses any property of
the innovation law beyond $(0,\ivar)$, the same estimator is obtained for
every innovation distribution with those moments.
\end{proof}

Theorem~\ref{thm:l-lmmse} gives a precise interpretation of the approximation:
Gaussian-projected BP is not merely a heuristic Gaussianisation of the
messages; it computes the \emph{best second-order linear estimator}. One
asymmetry in the statement deserves emphasis: the \emph{estimator} is the
same for every innovation law with matched first two moments, but its
\emph{error} against the true posterior mean is not --- each law has its own
nonlinear posterior, and the sweep of Section~\ref{sec:l-innovation-sweep}
measures exactly how differently the same estimator misses different truths.
Comparing it with grid BP therefore measures the loss from discarding
non-Gaussian structure beyond the first two moments, law by law.

\subsection{Three metrics, and which one is primary}
\label{sec:l-metrics}

The obvious quantity to report is the relative score error, and it is what the
diffusion model ultimately consumes. But by \eqref{eq:bp-master-error} it is
the posterior-mean error multiplied by $e^{-t}/\Deltat$ and divided by the norm
of the reference score, and both of those move by orders of magnitude across
the schedule. A score error that falls with $t$ therefore cannot, on its own,
distinguish an approximation that is becoming more accurate from one whose
errors are simply being scaled away. The two effects have to be separated, and
the posterior mean is where they separate.

\paragraph{The primary metric.} The exact posterior mean $m(x) = \E[a \mid x]$
is the MMSE estimator, so the distance to it has a risk interpretation that no
other choice has.

\begin{proposition}[Excess-risk decomposition]\label{prop:l-excess}
Let $m(x) = \E[a\mid x]$ and let $\widehat a(x)$ be any estimator with finite
second moment. Then
\begin{equation}
  \underbrace{\E\big\|a - \widehat a(x)\big\|^2}_{R(\widehat a)}
  \;=\;
  \underbrace{\E\big\|a - m(x)\big\|^2}_{R_{\mathrm{MMSE}}}
  \;+\;
  \E\big\|m(x) - \widehat a(x)\big\|^2 ,
  \label{eq:l-excess}
\end{equation}
so the squared distance between the two posterior means is \emph{exactly} the
excess denoising risk incurred by using $\widehat a$ in place of the exact
posterior mean.
\end{proposition}

\begin{proof}
Write $a - \widehat a = (a - m) + (m - \widehat a)$ and expand. The cross term
is $2\,\E[(a-m)^\T(m-\widehat a)]$; conditioning on $x$ and using
$\E[a - m(x) \mid x] = 0$, together with the fact that both $m(x)$ and
$\widehat a(x)$ are functions of $x$ alone, kills it. \end{proof}

Applied with $\widehat a = \widehat a_{\mathrm{LMMSE}}$, which
Theorem~\ref{thm:l-lmmse} identifies as what the closure returns,
\eqref{eq:l-excess} says that the gap between the two posterior means is not
an arbitrary distance but the price in mean squared error of denoising with
second moments only. We report it normalised by the signal energy, which for
this chain is $\E\|a\|^2 = \nsites$:
\begin{equation}
  \varepsilon_{\mathrm{PM}}
  \;=\;
  \sqrt{\frac{1}{\nsites}\,
        \E\big\| \widehat a_{\mathrm{LMMSE}} - \E[a \mid x] \big\|^2 } ,
  \label{eq:l-epspm}
\end{equation}
a dimensionless number, free of the prefactor, and stable at both ends
of the schedule. Normalising by $\E\|m(x)\|^2$ instead would be equally
natural at small $t$ but degenerates as $t\to\infty$, where the posterior mean
itself vanishes.

\paragraph{The secondary metrics.} Two more are reported alongside. The
\emph{excess risk fraction} $\E\|m - \widehat a\|^2 / R_{\mathrm{MMSE}}$ turns
\eqref{eq:l-excess} into a statement a practitioner can act on --- the
percentage by which closure raises the denoising error above the floor set by
the posterior itself --- and, because the floor grows with $t$ while the excess
shrinks, it is the sharpest of the three. And the relative score error is
retained, because it is what \eqref{eq:bp-master-error} converts the inference
error into, and therefore what the generative process actually sees.

\subsection{The measurement}

\begin{center}
\begin{minipage}{\linewidth}\centering
\captionof{table}[Gaussian-closure error across diffusion time]{Gaussian
closure against validated grid BP on the Laplace chain ($\alpha=0.85$,
$\nsites=32$, $60$ trials per $t$, reference $\ngrid=801$); all columns from
one sweep, sampled densely for $t\lesssim0.3$. The left group measures the
posterior approximation directly, before the  prefactor is applied:
$\varepsilon_{\mathrm{PM}}$ \eqref{eq:l-epspm} with its $10$--$90\%$ band, the
same normalised to $t=0.02$, and the excess denoising MSE of
Proposition~\ref{prop:l-excess} as a percentage of the MMSE floor (also
listed, since it grows while the excess shrinks). The last column is the
resulting discrepancy in the diffusion score.}
\label{tab:l-closure}
\small
\setlength{\tabcolsep}{5pt}
\renewcommand{\arraystretch}{1.1}
\begin{tabular}{@{}lccccc c@{}}
\toprule
& \multicolumn{5}{c}{\textit{posterior inference}}
& \textit{score} \\
\cmidrule(lr){2-6}\cmidrule(l){7-7}
$t$
& $\varepsilon_{\mathrm{PM}}$
& $10$--$90\%$
& $\varepsilon_{\mathrm{PM}}/\varepsilon_{\mathrm{PM}}^{\,0.02}$
& \shortstack{excess MSE\\{}[\% floor]}
& \shortstack{MMSE\\floor}
& \shortstack{rel.\ score\\error} \\
\midrule
$0.02$  & $0.036$  & $(0.029,\,0.052)$   & $1.00$ & $3.9$   & $0.032$ & $0.38$ \\
$0.032$ & $0.040$  & $(0.034,\,0.057)$   & $1.14$ & $3.4$   & $0.047$ & $0.30$ \\
$0.05$  & $0.050$  & $(0.036,\,0.079)$   & $1.42$ & $3.9$   & $0.066$ & $0.27$ \\
$0.072$ & $0.053$  & $(0.038,\,0.081)$   & $\mathbf{1.50}$ & $3.1$ & $0.089$ & $0.20$ \\
$0.105$ & $0.050$  & $(0.037,\,0.100)$   & $1.41$ & $2.1$   & $0.114$ & $0.14$ \\
$0.153$ & $0.043$  & $(0.034,\,0.085)$   & $1.20$ & $1.2$   & $0.151$ & $0.10$ \\
$0.222$ & $0.038$  & $(0.033,\,0.058)$   & $1.06$ & $0.74$  & $0.193$ & $0.065$ \\
$0.322$ & $0.033$  & $(0.027,\,0.059)$   & $0.93$ & $0.43$  & $0.253$ & $0.041$ \\
\addlinespace
$0.98$  & $0.016$  & $(0.012,\,0.028)$   & $0.45$ & $0.045$ & $0.584$ & $7.3\times10^{-3}$ \\
$2.07$  & $0.0018$ & $(0.0013,\,0.0038)$ & $0.05$ & $<0.01$ & $0.916$ & $2.4\times10^{-4}$ \\
\bottomrule
\end{tabular}
\end{minipage}
\end{center}

\paragraph{Two readings that disagree.}
Read the last column and the closure improves steadily as noise is added: the
relative score error falls by a factor of $3.6$ between $t=0.02$ and $t=0.15$.
Read $\varepsilon_{\mathrm{PM}}$ and it does not. The posterior-mean error
\emph{rises} by half over $t\in[0.02,0.07]$, is still larger at $t=0.15$ than
at $t=0.02$, and begins to decay only past $t\simeq0.3$
(Figure~\ref{fig:l-metrics}, left and centre). By \eqref{eq:bp-master-error}
the gap between the two readings is the prefactor $e^{-t}/\Deltat\sim1/2t$ and
the shrinking norm of the reference score, neither of which reports on the
quality of the inference. Relative score error measures the \emph{consequence}
of the approximation for the diffusion; $\varepsilon_{\mathrm{PM}}$ measures
the inference itself, and only the latter is flat at small $t$.

\paragraph{As a denoiser.}
By Proposition~\ref{prop:l-excess} the squared posterior-mean discrepancy is
exactly the excess MSE over the MMSE estimator. As a fraction of the MMSE
floor it falls from $3.9\%$ at $t=0.05$ to $0.43\%$ at $t=0.32$ and below
$0.05\%$ once $t\gtrsim1$, while the floor itself grows from $0.07$ to $0.92$
over the same range: the approximation error shrinks as the irreducible
denoising uncertainty grows, and the ratio is the sharpest of the metrics for
that reason.

\paragraph{Direction survives magnitude.}
The probability of a trial exceeding $10\%$ relative score error is $1.00$
through $t\simeq0.07$, passes $0.55$ at $t=0.15$, and reaches zero by
$t\simeq0.7$. Direction is far more robust: the median cosine similarity to
the reference score is $0.964$ at $t=0.05$, exceeds $0.99$ by $t\simeq0.11$,
and is numerically one near $t=1$ (Figure~\ref{fig:bp-closure}, right).
Gaussian closure loses \emph{how strongly} to move before it loses
\emph{where} to move.

\paragraph{Why the peak is where it is.}
The maximum in $\varepsilon_{\mathrm{PM}}$ separates two regimes. As $t\to0$
each latent state is tightly pinned by its own near-noiseless observation, so
inference is easy even for a non-Gaussian prior. At large $t$ Gaussian noise
suppresses the higher-order structure of the innovation law and the tilted
posteriors approach Gaussianity. The largest error therefore sits between the
limits, where the evidence is weak enough for the prior's shape to matter and
strong enough that it has not yet been washed out. In the worst individual
trials the reference posterior is strongly non-Gaussian --- heavy-tailed and
sharply kinked, though never bimodal: the negative log-posterior of the
Laplace chain under Gaussian observations is strictly convex, so the posterior
and its marginals are log-concave and unimodal at every $t > 0$ --- whereas a
single Gaussian retains only one mean and variance. (Genuinely bimodal
posteriors do arise for the bimodal mixture innovations of
Section~\ref{sec:l-innovation-sweep}, and they are that sweep's hardest case.)
This is a qualitative mechanism for the aggregate curve, not a direct
measurement of posterior non-Gaussianity.

\begin{figure}[t]
\centering
\includegraphics[width=\linewidth]{fig_laplace_metrics.pdf}
\caption[The Gaussian closure under three metrics]{
The same experiment under three metrics ($\alpha=0.85$, $\nsites=32$, $60$
trials per $t$, reference $\ngrid=801$).
\emph{Left:} the posterior-mean error \eqref{eq:l-epspm}, median and
$10$--$90\%$ band. It is not monotone: it peaks at $t\simeq0.07$ and is no
smaller at $t=0.15$ than at $t=0.02$.
\emph{Centre:} the same runs seen through the relative score error, which
falls by $3.6\times$ over exactly the interval on which
$\varepsilon_{\mathrm{PM}}$ is flat. The gap between the curves is the
 prefactor of \eqref{eq:bp-master-error} together with the changing
norm of the reference score; the two cross near $t\simeq0.5$.
\emph{Right:} the excess denoising MSE of Proposition~\ref{prop:l-excess}, as
a percentage of the MMSE floor.}
\label{fig:l-metrics}
\end{figure}

\begin{figure}[t]
\centering
\includegraphics[width=\linewidth]{fig_laplace_closure.pdf}
\caption[Gaussian-closure score error against grid BP]{
The score-side view of the same sweep.
\emph{Left:} median relative score error with its $10$--$90\%$ band, against
the grid-BP self-convergence error ($\ngrid=401$ versus $801$) on identical
data. The closure error exceeds the numerical reference budget by at least
three orders of magnitude everywhere and by four to five over most of the
range, so the measurement is limited by the closure, not the quadrature.
\emph{Right:} relative score error against $1-\cos$ of the angle between the
projected and reference score vectors. Directional error is smaller than
magnitude error at every diffusion time, by a factor of about five at
$t=0.02$ and by three to five orders of magnitude once $t\gtrsim1$.}
\label{fig:bp-closure}
\end{figure}

\subsection{Varying the innovation law}\label{sec:l-innovation-sweep}

Holding $\operatorname{Cov}(a_i,a_j)=\alpha^{|i-j|}$ fixed, we vary the
innovation distribution: Laplace, Student-$t$ with $\nu\in\{5,8\}$, and
symmetric Gaussian mixtures of increasing separation. Across every family
substantial closure error is confined to $t\lesssim0.4$ and falls to about
$10^{-3}$ or below by $t=2.4$. At $t=0.05$ and $\alpha=0.5$:

\begin{center}
\small
\begin{tabular}{@{}lc@{}}
\toprule
innovation law & median rel.\ score error \\
\midrule
strongly bimodal Gaussian mixture & $0.94$ \\
Laplace                           & $0.52$ \\
Student-$t$, $\nu=5$              & $0.32$ \\
mildly bimodal Gaussian mixture   & $0.07$ \\
\bottomrule
\end{tabular}
\end{center}

\noindent Strong bimodality is the hardest case: a single Gaussian cannot
place mass on two separated posterior modes and sits between them. Stronger
chain correlation helps --- for Laplace innovations at $t=0.05$ the median
error falls from roughly one half at $\alpha=0.5$ to $0.27$ at $\alpha=0.85$
and $0.13$ at $\alpha=0.95$ --- so informative neighbours reduce the weight
the local non-Gaussian shape carries, again an interpretation rather than a
direct measurement of posterior Gaussianity.